\chapter{ログランク検定}\label{ch:m02}

\begin{priobox}{A}
\kakomon{2016}{3}〔2〕では手順が文章で説明されるが，「リスク集合の数え方」「期待イベント数 $n_{1j}/n_j$」「分散 $\frac{n_{1j}}{n_j}\frac{n_{0j}}{n_j}$」「$\chi^2=U^2/V$ が自由度1」を知らないと時間内に計算を終えられない．
群逐次デザイン（\kakomon{2017}{2}，第9章）でも検定統計量として現れる．本編の第9章9.5節．
\end{priobox}

% ============================================================
\section{概要}\label{sec:m02-gaiyo}
2群の生存関数が等しい（$H_0:S_1(t)=S_0(t)$，すべての $t$）かを検定する．
各イベント時点で「その直前に観察中の人（リスク集合）」を群×（イベントあり・なし）の $2\times2$ 表にまとめ，
群1の観測イベント数 $O$ と，$H_0$ の下での期待イベント数 $E$ の差を全時点で足す．

\begin{pointbox}[覚える式]
時点 $t_{(j)}$ のリスク集合の人数を $n_{1j},n_{0j}$（計 $n_j$），イベント数を $d_{1j}$（計 $d_j$）とすると
\[
E_{1j}=d_j\frac{n_{1j}}{n_j},\qquad
V_j=\frac{d_j\,n_{1j}\,n_{0j}\,(n_j-d_j)}{n_j^2(n_j-1)}\quad\Bigl(d_j=1\text{ なら }V_j=\frac{n_{1j}}{n_j}\cdot\frac{n_{0j}}{n_j}\Bigr),
\]
\[
U=\sum_j(d_{1j}-E_{1j})=O_1-E_1,\qquad V=\sum_jV_j,\qquad \chi^2_{\mathrm{LR}}=\frac{U^2}{V}\ \xrightarrow{d}\ \chi^2_1\quad(H_0\text{ の下}).
\]
$P$ 値は $\sqrt{\chi^2}$ を標準正規分布表で引いて上側確率を2倍する．
\end{pointbox}

% ============================================================
\section{詳細}\label{sec:m02-shosai}
\subsection{表の作り方}
\begin{enumerate}
\item 両群をまとめて観察時間を昇順に並べる．
\item \textbf{イベントが起きた時点だけ}を取り出す（打ち切りだけの時点は行を作らない）．
\item 各イベント時点 $t_{(j)}$ で，$t_{(j)}$ 以上の観察時間をもつ人を群ごとに数えて $n_{1j},n_{0j}$ とする．
      同じ時点にイベントと打ち切りがあれば\textbf{イベントが先}（その時点の打ち切り例はリスク集合に含める）．
\item $E_{1j}$，$V_j$ を計算して足す．一方の群のリスク集合が0人になった時点では $E_{1j}=d_j$ または0，$V_j=0$ で，検定に寄与しない．
\end{enumerate}

\subsection{期待値と分散の導出}
時点 $t_{(j)}$ の表で周辺度数（$n_{1j},n_{0j},d_j$）を固定する．
$H_0$ の下では，$n_j$ 人のうち誰がイベントを起こすかは群によらないので，
$d_{1j}$ は「$n_j$ 人（うち群1が $n_{1j}$ 人）から $d_j$ 人を非復元で選んだときに含まれる群1の人数」であり，超幾何分布に従う．
超幾何分布の平均と分散から
\[
\E[d_{1j}]=d_j\frac{n_{1j}}{n_j},\qquad \V[d_{1j}]=d_j\frac{n_{1j}}{n_j}\Bigl(1-\frac{n_{1j}}{n_j}\Bigr)\frac{n_j-d_j}{n_j-1}
=\frac{d_jn_{1j}n_{0j}(n_j-d_j)}{n_j^2(n_j-1)} .
\]
最後の因子 $(n_j-d_j)/(n_j-1)$ は非復元抽出の有限母集団修正である．
$d_j=1$ ならこの因子は1で，$V_j$ は「確率 $n_{1j}/n_j$ の Bernoulli 分布の分散」になる（\kakomon{2016}{3} の問題文の「二項分布の分散」はこれ）．
\paragraph{同時イベントの小さな例}
$n_{1j}=3$，$n_{0j}=4$，$d_j=2$ なら $E_{1j}=2\times3/7=0.857$，$V_j=\dfrac{2\times3\times4\times5}{7^2\times6}=\dfrac{120}{294}=0.408$．

\subsection{時点ごとの分散を足してよい理由}
時点ごとの表は独立ではない（前の時点のイベントが後のリスク集合を決める）．
しかし $H_0$ の下で，時点 $t_{(j)}$ の直前までの情報（リスク集合の大きさを含む）と $d_j$ を与えると $d_{1j}-E_{1j}$ の条件付き期待値は0である．
$i<j$ なら $d_{1i}-E_{1i}$ はその情報で決まるので，反復期待値により
\[
\E[(d_{1i}-E_{1i})(d_{1j}-E_{1j})]=\E\bigl[(d_{1i}-E_{1i})\,\E[d_{1j}-E_{1j}\mid\text{過去},d_j]\bigr]=0 .
\]
時点間の共分散が0なので $\V[U]=\sum_j\E[V_j]$ で，$V=\sum_jV_j$ はその推定量である．$\chi^2_1$ への収束はマルチンゲールの中心極限定理による（証明は略）．

\subsection{性質}
\begin{itemize}
\item \textbf{Cox モデルのスコア検定と一致}：$h(t\mid x)=h_0(t)e^{\beta x}$（$x$ は群の指示）の $H_0:\beta=0$ のスコア検定はログランク検定になる（第3章）．
      したがって比例ハザード型の差（ハザード比が一定で1に近い）に対して最も検出力が高い．
\item \textbf{生存曲線が交差すると検出力を失う}：前半と後半で $O-E$ の符号が逆になり打ち消し合う．
\item どちらの群で計算しても $\chi^2$ は同じ（両群の $O-E$ の和は0）．
\item 層別 Cochran（CMH）検定（第7章）と同じ形：時点を層とみなしたもの．
\end{itemize}

\begin{warnbox}
\begin{itemize}
\item 打ち切りだけの時点にも行を作る．同時点の打ち切りを先に処理してリスク集合を1減らす．
\item 分散を $\sum E_{1j}$（ポアソン近似）や $\sum E_{1j}+\sum E_{0j}$ で代用する．問題文の指定（二項分散・超幾何分散）に従う．
\item $\chi^2$ を自由度2で判定する（2群なら自由度1）．
\item 有意でないことを「生存曲線が同じ」の根拠にする．とくに曲線が交差している場合．
\end{itemize}
\end{warnbox}

% ============================================================
\section{出題例}\label{sec:m02-shutsudai}
\begin{shutsudai}{\kakomon{2016}{3}〔2〕：2群各5例のログランク検定}
イベント時点ごとのイベント数とリスク集合の人数の表（表2）が与えられ，各時点のイベントは1例．
スコア $u=\sum_t(d_t-n_{1t}/n_t)$ と分散 $v=\sum_t\frac{n_{1t}}{n_t}\frac{n_t-n_{1t}}{n_t}$ の作り方は文章で説明される．
群1で $u=-2.3222$，$v=0.9591$，$\chi^2=u^2/v=5.62$，$\sqrt{5.62}=2.37$ から両側 $P=2\times0.0089=0.018$．\\
この問題では群1の最後の2時点で群2のリスク集合が0人になり，$u,v$ への寄与が0になる．
\end{shutsudai}
\begin{shutsudai}{\kakomon{2017}{2}：群逐次デザインの検定統計量}
中間解析・最終解析の検定統計量としてログランク検定の $Z(t)$ が使われる．$Z(t)$ が $H_0$ の下で標準正規分布に従うことは与えられる．
情報量（$V$）がイベント数にほぼ比例するので，情報時間を（中間解析時のイベント数）/（最終の目標イベント数）とする（第9章）．
\end{shutsudai}

% ============================================================
\section{例題}\label{sec:m02-reidai}
\begin{reidai}[\kakomon{2016}{3}〔2〕型：ログランク検定の手計算]\label{reidai:m02-1}
2群各5人の生存時間データを次に示す（打ち切りの有無は 1：打ち切り，0：イベント）．
\begin{center}\small
\begin{tabular}{@{}cccc@{\qquad}cccc@{}}
\toprule
群 & 観察時間（週） & 打ち切り & & 群 & 観察時間（週） & 打ち切り\\
\midrule
1 & 3 & 1 & & 2 & 1 & 0\\
1 & 6 & 0 & & 2 & 2 & 0\\
1 & 8 & 1 & & 2 & 4 & 0\\
1 & 10 & 0 & & 2 & 5 & 1\\
1 & 12 & 0 & & 2 & 7 & 0\\
\bottomrule
\end{tabular}
\end{center}
\begin{enumerate}[label=〔\arabic*〕]
\item イベント時点ごとに，各群のイベント数とリスク集合の人数の表を作れ．
\item 群1について，各イベント時点の期待イベント数と分散（各時点のイベントは1例なので Bernoulli 分散）を求め，ログランク検定の $\chi^2$ 統計量と両側 $P$ 値を求めよ．
\item 結果を解釈せよ．
\end{enumerate}
\end{reidai}

% ============================================================
\section{例題の詳しい解説}\label{sec:m02-kaisetsu}
\begin{kaitou}{reidai:m02-1}
〔1〕打ち切りの列は「1：打ち切り」なので，群1のイベントは6，10，12週，打ち切りは3，8週．群2のイベントは1，2，4，7週，打ち切りは5週．
全体を並べると $1,\ 2,\ 3^+,\ 4,\ 5^+,\ 6,\ 7,\ 8^+,\ 10,\ 12$．イベント時点は 1，2，4，6，7，10，12 の7つ．
各時点で，その時点以上の観察時間をもつ人を数える．
\begin{center}\small
\begin{tabular}{@{}lccccccc@{}}
\toprule
観察時間（週） & 1 & 2 & 4 & 6 & 7 & 10 & 12\\
\midrule
群1のイベント数 $d_{1j}$ & 0 & 0 & 0 & 1 & 0 & 1 & 1\\
群2のイベント数 & 1 & 1 & 1 & 0 & 1 & 0 & 0\\
群1のリスク集合 $n_{1j}$ & 5 & 5 & 4 & 4 & 3 & 2 & 1\\
群2のリスク集合 $n_{0j}$ & 5 & 4 & 3 & 1 & 1 & 0 & 0\\
リスク集合の総数 $n_j$ & 10 & 9 & 7 & 5 & 4 & 2 & 1\\
\bottomrule
\end{tabular}
\end{center}
（4週の群1は3週の打ち切りで4人．6週の群2は，1・2・4週のイベントと5週の打ち切りで1人（7週の人）．
7週の群1は6週のイベントで3人．10週の群1は8週の打ち切りで2人．）

〔2〕各時点のイベントは1例なので $E_{1j}=n_{1j}/n_j$，$V_j=n_{1j}n_{0j}/n_j^2$．
\begin{center}\small
\begin{tabular}{@{}lcccccccc@{}}
\toprule
観察時間 & 1 & 2 & 4 & 6 & 7 & 10 & 12 & 計\\
\midrule
$d_{1j}$ & 0 & 0 & 0 & 1 & 0 & 1 & 1 & $O_1=3$\\
$E_{1j}=n_{1j}/n_j$ & $5/10$ & $5/9$ & $4/7$ & $4/5$ & $3/4$ & $2/2$ & $1/1$ & \\
 & 0.5000 & 0.5556 & 0.5714 & 0.8000 & 0.7500 & 1 & 1 & $E_1=5.1770$\\
$d_{1j}-E_{1j}$ & $-0.5000$ & $-0.5556$ & $-0.5714$ & $0.2000$ & $-0.7500$ & 0 & 0 & $U=-2.1770$\\
$V_j=n_{1j}n_{0j}/n_j^2$ & $25/100$ & $20/81$ & $12/49$ & $4/25$ & $3/16$ & 0 & 0 & \\
 & 0.2500 & 0.2469 & 0.2449 & 0.1600 & 0.1875 & 0 & 0 & $V=1.0893$\\
\bottomrule
\end{tabular}
\end{center}
\[
\chi^2_{\mathrm{LR}}=\frac{U^2}{V}=\frac{(-2.1770)^2}{1.0893}=\frac{4.7393}{1.0893}=4.35 .
\]
$\sqrt{4.35}=2.09$ で，正規分布表より $\Prob(Z>2.09)=0.0183$，両側 $P=2\times0.0183=0.037$．
（$\chi^2$ 分布表の $\chi^2_1(0.05)=3.841<4.35$ と比べてもよい．）有意水準5\%で $H_0$ を棄却する．

〔3〕群1の観測イベント数3は期待値5.18より少なく，群1のほうが生存が良い（イベントが遅い）方向の差が有意である．
ただし各群5人の小標本で，$\chi^2_1$ は漸近近似である．差の大きさ（ハザード比）はログランク検定からはわからないので，Cox 回帰（第3章）で推定する．
10週と12週は群2のリスク集合が0人なので $U,V$ に寄与しない（群2が全員いなくなった後のイベントは比較の情報を持たない）．
\end{kaitou}

\begin{summarybox}
\begin{itemize}
\item イベント時点ごとの $2\times2$ 表．同時点ではイベントが先．打ち切りだけの時点は行を作らない．
\item $E_{1j}=d_jn_{1j}/n_j$，$V_j=d_jn_{1j}n_{0j}(n_j-d_j)/\{n_j^2(n_j-1)\}$（$d_j=1$ なら $n_{1j}n_{0j}/n_j^2$）．
\item $\chi^2=(O_1-E_1)^2/V$，自由度1．Cox のスコア検定と一致し，比例ハザード型の差に強く，交差に弱い．
\end{itemize}
\end{summarybox}
